\begin{table}[H]
\centering
\caption{Matriz cualitativa de las alternativas de diseño. La decisión final depende del propósito de uso y de los resultados del protocolo común; no se asignaron pesos retrospectivos.}
\label{tab:chapter0-decision-matrix}
\footnotesize
\setlength{\tabcolsep}{3pt}
\renewcommand{\arraystretch}{1.16}
\resizebox{\textwidth}{!}{%
\begin{tabular}{p{0.14\textwidth} p{0.21\textwidth} p{0.20\textwidth} p{0.24\textwidth} p{0.13\textwidth}}
\hline
Alternativa & Forma de incorporar contexto & Salida que puede inspeccionarse & Restricción relevante & Papel en la comparación \tabularnewline
\hline
HMM finito & Transiciones entre un número fijo de estados latentes & Trayectoria de estados y probabilidades de transición & La capacidad depende de la rejilla elegida y el costo crece al ampliar el número de estados & Referencia latente de capacidad controlada \tabularnewline
HDP-HMM truncado & Inventario jerárquico de estados, limitado por una cota de cómputo & Trayectoria y número efectivo de estados ocupados & Requiere revisar convergencia, deriva y efecto de la truncación & Alternativa latente flexible \tabularnewline
VOMM & Contextos observables de longitud variable & Conteos y probabilidades por contexto & La longitud máxima se selecciona en validación y no hay estados latentes & Referencia predictiva explícita \tabularnewline
Transformador causal & Atención sobre una ventana de eventos anteriores & Distribución predictiva del siguiente evento & Exige más parámetros y su atención no equivale a una trayectoria discreta & Alternativa de mayor capacidad contextual \tabularnewline
\hline
\end{tabular}
}
\end{table}
